\makeatletter
\newlabel{ej:parciales}{{1.2}{1}{}{ejemplo.2}{}}
\newlabel{ej:parciales@cref}{{[ejemplo][2][]1.2}{[1][1][]1}}
\newlabel{ej:direccional}{{1.10}{2}{}{ejemplo.10}{}}
\newlabel{ej:direccional@cref}{{[ejemplo][10][]1.10}{[1][2][]2}}
\newlabel{fig:nivel}{{1.1}{4}{Curvas de nivel $x^{2}+3y^{2}=c$ para $c=1$; $2{,}5$; $4$. En $(1,1)$ el gradiente (flecha, a escala reducida) es perpendicular a la tangente a la curva $f=4$ (trazo discontinuo), de dirección $(3,-1)\T $}{figure.1.1}{}}
\newlabel{fig:nivel@cref}{{[figure][1][1]1.1}{[1][3][]4}}
\newlabel{ej:dg}{{1.16}{4}{}{ejemplo.16}{}}
\newlabel{ej:dg@cref}{{[ejemplo][16][]1.16}{[1][4][]4}}
\newlabel{eje:1-vf}{{1.1}{5}{}{ejercicio.1.1}{}}
\newlabel{eje:1-vf@cref}{{[ejercicio][1][1]1.1}{[1][4][]5}}
\newlabel{eje:1-gradiente}{{1.2}{5}{}{ejercicio.1.2}{}}
\newlabel{eje:1-gradiente@cref}{{[ejercicio][2][1]1.2}{[1][5][]5}}
\newlabel{eje:1-d2l}{{1.3}{5}{}{ejercicio.1.3}{}}
\newlabel{eje:1-d2l@cref}{{[ejercicio][3][1]1.3}{[1][5][]5}}
\newlabel{eje:1-norma}{{1.4}{5}{}{ejercicio.1.4}{}}
\newlabel{eje:1-norma@cref}{{[ejercicio][4][1]1.4}{[1][5][]5}}
\newlabel{eje:1-mit}{{1.5}{5}{}{ejercicio.1.5}{}}
\newlabel{eje:1-mit@cref}{{[ejercicio][5][1]1.5}{[1][5][]5}}
\newlabel{eje:1-minimo}{{1.6}{5}{}{ejercicio.1.6}{}}
\newlabel{eje:1-minimo@cref}{{[ejercicio][6][1]1.6}{[1][5][]5}}
\newlabel{ej:jac-lineal}{{2.2}{6}{}{ejemplo.2}{}}
\newlabel{ej:jac-lineal@cref}{{[ejemplo][2][]2.2}{[1][6][]6}}
\newlabel{ej:jac}{{2.3}{6}{}{ejemplo.3}{}}
\newlabel{ej:jac@cref}{{[ejemplo][3][]2.3}{[1][6][]6}}
\newlabel{ej:cadena}{{2.6}{7}{}{ejemplo.6}{}}
\newlabel{ej:cadena@cref}{{[ejemplo][6][]2.6}{[1][7][]7}}
\newlabel{fig:grafo}{{2.1}{8}{Grafo computacional del \cref {ej:cadena}. Dentro de cada nodo figura el valor calculado en el paso hacia delante; junto a él, en gris, el adjunto calculado en el paso hacia atrás con la \cref {prop:retro}. La entrada $x$ recibe contribuciones por dos caminos, que se suman}{figure.2.1}{}}
\newlabel{fig:grafo@cref}{{[figure][1][2]2.1}{[1][7][]8}}
\newlabel{eje:2-formas}{{2.1}{9}{}{ejercicio.2.1}{}}
\newlabel{eje:2-formas@cref}{{[ejercicio][1][2]2.1}{[1][9][]9}}
\newlabel{eje:2-producto}{{2.2}{9}{}{ejercicio.2.2}{}}
\newlabel{eje:2-producto@cref}{{[ejercicio][2][2]2.2}{[1][9][]9}}
\newlabel{eje:2-jacobiano}{{2.3}{9}{}{ejercicio.2.3}{}}
\newlabel{eje:2-jacobiano@cref}{{[ejercicio][3][2]2.3}{[1][9][]9}}
\newlabel{eje:2-d2l}{{2.4}{9}{}{ejercicio.2.4}{}}
\newlabel{eje:2-d2l@cref}{{[ejercicio][4][2]2.4}{[1][9][]9}}
\newlabel{eje:2-coste}{{2.5}{9}{}{ejercicio.2.5}{}}
\newlabel{eje:2-coste@cref}{{[ejercicio][5][2]2.5}{[1][9][]9}}
\newlabel{eje:2-escalar}{{2.6}{9}{}{ejercicio.2.6}{}}
\newlabel{eje:2-escalar@cref}{{[ejercicio][6][2]2.6}{[1][9][]9}}
\newlabel{ej:capa}{{3.4}{11}{}{ejemplo.4}{}}
\newlabel{ej:capa@cref}{{[ejemplo][4][]3.4}{[1][11][]11}}
\newlabel{ej:p-y}{{3.9}{12}{}{ejemplo.9}{}}
\newlabel{ej:p-y@cref}{{[ejemplo][9][]3.9}{[1][12][]12}}
\newlabel{eje:3-embeddings}{{3.1}{13}{}{ejercicio.3.1}{}}
\newlabel{eje:3-embeddings@cref}{{[ejercicio][1][3]3.1}{[1][13][]13}}
\newlabel{eje:3-signos}{{3.2}{13}{}{ejercicio.3.2}{}}
\newlabel{eje:3-signos@cref}{{[ejercicio][2][3]3.2}{[1][13][]13}}
\newlabel{eje:3-capa}{{3.3}{13}{}{ejercicio.3.3}{}}
\newlabel{eje:3-capa@cref}{{[ejercicio][3][3]3.3}{[1][13][]13}}
\newlabel{eje:3-softmax}{{3.4}{13}{}{ejercicio.3.4}{}}
\newlabel{eje:3-softmax@cref}{{[ejercicio][4][3]3.4}{[1][13][]13}}
\newlabel{eje:3-nucleo}{{3.5}{13}{}{ejercicio.3.5}{}}
\newlabel{eje:3-nucleo@cref}{{[ejercicio][5][3]3.5}{[1][13][]13}}
\newlabel{eje:3-word2vec}{{3.6}{13}{}{ejercicio.3.6}{}}
\newlabel{eje:3-word2vec@cref}{{[ejercicio][6][3]3.6}{[1][13][]13}}
\newlabel{ej:hessiano}{{4.3}{14}{}{ejemplo.3}{}}
\newlabel{ej:hessiano@cref}{{[ejemplo][3][]4.3}{[1][14][]14}}
\newlabel{ej:taylor}{{4.6}{15}{}{ejemplo.6}{}}
\newlabel{ej:taylor@cref}{{[ejemplo][6][]4.6}{[1][15][]15}}
\newlabel{ej:dg-cuad}{{4.13}{16}{}{ejemplo.13}{}}
\newlabel{ej:dg-cuad@cref}{{[ejemplo][13][]4.13}{[1][16][]16}}
\newlabel{eje:4-vf}{{4.1}{17}{}{ejercicio.4.1}{}}
\newlabel{eje:4-vf@cref}{{[ejercicio][1][4]4.1}{[1][17][]17}}
\newlabel{eje:4-taylor}{{4.2}{17}{}{ejercicio.4.2}{}}
\newlabel{eje:4-taylor@cref}{{[ejercicio][2][4]4.2}{[1][17][]17}}
\newlabel{eje:4-mit}{{4.3}{18}{}{ejercicio.4.3}{}}
\newlabel{eje:4-mit@cref}{{[ejercicio][3][4]4.3}{[1][18][]18}}
\newlabel{eje:4-tasa}{{4.4}{18}{}{ejercicio.4.4}{}}
\newlabel{eje:4-tasa@cref}{{[ejercicio][4][4]4.4}{[1][18][]18}}
\newlabel{eje:4-necesaria}{{4.5}{18}{}{ejercicio.4.5}{}}
\newlabel{eje:4-necesaria@cref}{{[ejercicio][5][4]4.5}{[1][18][]18}}
\newlabel{fig:convexa}{{5.1}{19}{Una función convexa queda por debajo de sus cuerdas (\cref {def:convexa}) y por encima de sus tangentes (\cref {teo:primer-orden})}{figure.5.1}{}}
\newlabel{fig:convexa@cref}{{[figure][1][5]5.1}{[1][19][]19}}
\newlabel{ej:cuad-convexa}{{5.6}{20}{}{ejemplo.6}{}}
\newlabel{ej:cuad-convexa@cref}{{[ejemplo][6][]5.6}{[1][20][]20}}
\newlabel{ej:red-lineal}{{5.13}{22}{Una red lineal de dos capas}{ejemplo.13}{}}
\newlabel{ej:red-lineal@cref}{{[ejemplo][13][]5.13}{[1][22][]22}}
\newlabel{fig:hiperbola}{{5.2}{23}{Mínimos globales de $L(w_1,w_2)=(1-w_1w_2)^{2}$ (\cref {ej:red-lineal}). El segmento que une dos mínimos sale del conjunto de mínimos: la pérdida no es convexa}{figure.5.2}{}}
\newlabel{fig:hiperbola@cref}{{[figure][2][5]5.2}{[1][22][]23}}
\newlabel{eje:5-vf}{{5.1}{23}{}{ejercicio.5.1}{}}
\newlabel{eje:5-vf@cref}{{[ejercicio][1][5]5.1}{[1][23][]23}}
\newlabel{eje:5-hessianos}{{5.2}{23}{}{ejercicio.5.2}{}}
\newlabel{eje:5-hessianos@cref}{{[ejercicio][2][5]5.2}{[1][23][]23}}
\newlabel{eje:5-boyd}{{5.3}{23}{}{ejercicio.5.3}{}}
\newlabel{eje:5-boyd@cref}{{[ejercicio][3][5]5.3}{[1][23][]23}}
\newlabel{eje:5-red}{{5.4}{23}{}{ejercicio.5.4}{}}
\newlabel{eje:5-red@cref}{{[ejercicio][4][5]5.4}{[1][23][]23}}
\newlabel{eje:5-softmax-W}{{5.5}{23}{}{ejercicio.5.5}{}}
\newlabel{eje:5-softmax-W@cref}{{[ejercicio][5][5]5.5}{[1][23][]23}}
\newlabel{eje:5-unico}{{5.6}{23}{}{ejercicio.5.6}{}}
\newlabel{eje:5-unico@cref}{{[ejercicio][6][5]5.6}{[1][23][]23}}
\newlabel{ej:exp}{{6.3}{25}{}{ejemplo.3}{}}
\newlabel{ej:exp@cref}{{[ejemplo][3][]6.3}{[1][25][]25}}
\newlabel{fig:error-h}{{6.1}{27}{Error de la diferencia centrada para $f(x)=e^{x}$ en $x=0$, calculada en doble precisión (puntos), y cota $E(h)$ de la \cref {prop:h-optimo} con $\varepsilon =1{,}1\cdot 10^{-16}$ y $\abs {f}=M=1$ (línea). A la derecha domina el truncamiento (pendiente~2); a la izquierda, el redondeo (pendiente~$-1$). El error medido mínimo, $1{,}2\cdot 10^{-11}$, se obtiene con $h=10^{-5}$}{figure.6.1}{}}
\newlabel{fig:error-h@cref}{{[figure][1][6]6.1}{[1][26][]27}}
\newlabel{eje:6-porque}{{6.1}{28}{}{ejercicio.6.1}{}}
\newlabel{eje:6-porque@cref}{{[ejercicio][1][6]6.1}{[1][28][]28}}
\newlabel{eje:6-cubo}{{6.2}{28}{}{ejercicio.6.2}{}}
\newlabel{eje:6-cubo@cref}{{[ejercicio][2][6]6.2}{[1][28][]28}}
\newlabel{eje:6-direccional}{{6.3}{28}{}{ejercicio.6.3}{}}
\newlabel{eje:6-direccional@cref}{{[ejercicio][3][6]6.3}{[1][28][]28}}
\newlabel{eje:6-cs231n}{{6.4}{28}{}{ejercicio.6.4}{}}
\newlabel{eje:6-cs231n@cref}{{[ejercicio][4][6]6.4}{[1][28][]28}}
\newlabel{eje:6-fermi}{{6.5}{28}{}{ejercicio.6.5}{}}
\newlabel{eje:6-fermi@cref}{{[ejercicio][5][6]6.5}{[1][28][]28}}
\newlabel{eje:6-hoptimo}{{6.6}{28}{}{ejercicio.6.6}{}}
\newlabel{eje:6-hoptimo@cref}{{[ejercicio][6][6]6.6}{[1][28][]28}}
\newlabel{ap:soluciones}{{A}{29}{Soluciones de los ejercicios}{appendix.A}{}}
\newlabel{ap:soluciones@cref}{{[appendix][1][2147483647]A}{[1][29][]29}}
\newlabel{sol:eje:1-vf}{{A.1}{29}{Tema 1}{section*.3}{}}
\newlabel{sol:eje:1-vf@cref}{{[subappendix][1][2147483647,1]A.1}{[1][29][]29}}
\newlabel{sol:eje:1-gradiente}{{A.1}{29}{Tema 1}{section*.4}{}}
\newlabel{sol:eje:1-gradiente@cref}{{[subappendix][1][2147483647,1]A.1}{[1][29][]29}}
\newlabel{sol:eje:1-d2l}{{A.1}{30}{Tema 1}{section*.5}{}}
\newlabel{sol:eje:1-d2l@cref}{{[subappendix][1][2147483647,1]A.1}{[1][29][]30}}
\newlabel{sol:eje:1-norma}{{A.1}{30}{Tema 1}{section*.6}{}}
\newlabel{sol:eje:1-norma@cref}{{[subappendix][1][2147483647,1]A.1}{[1][30][]30}}
\newlabel{sol:eje:1-mit}{{A.1}{30}{Tema 1}{section*.7}{}}
\newlabel{sol:eje:1-mit@cref}{{[subappendix][1][2147483647,1]A.1}{[1][30][]30}}
\newlabel{sol:eje:1-minimo}{{A.1}{31}{Tema 1}{section*.8}{}}
\newlabel{sol:eje:1-minimo@cref}{{[subappendix][1][2147483647,1]A.1}{[1][31][]31}}
\newlabel{sol:eje:2-formas}{{A.2}{31}{Tema 2}{section*.9}{}}
\newlabel{sol:eje:2-formas@cref}{{[subappendix][2][2147483647,1]A.2}{[1][31][]31}}
\newlabel{sol:eje:2-producto}{{A.2}{32}{Tema 2}{section*.10}{}}
\newlabel{sol:eje:2-producto@cref}{{[subappendix][2][2147483647,1]A.2}{[1][32][]32}}
\newlabel{sol:eje:2-jacobiano}{{A.2}{32}{Tema 2}{section*.11}{}}
\newlabel{sol:eje:2-jacobiano@cref}{{[subappendix][2][2147483647,1]A.2}{[1][32][]32}}
\newlabel{sol:eje:2-d2l}{{A.2}{32}{Tema 2}{section*.12}{}}
\newlabel{sol:eje:2-d2l@cref}{{[subappendix][2][2147483647,1]A.2}{[1][32][]32}}
\newlabel{sol:eje:2-coste}{{A.2}{33}{Tema 2}{section*.13}{}}
\newlabel{sol:eje:2-coste@cref}{{[subappendix][2][2147483647,1]A.2}{[1][33][]33}}
\newlabel{sol:eje:2-escalar}{{A.2}{33}{Tema 2}{section*.14}{}}
\newlabel{sol:eje:2-escalar@cref}{{[subappendix][2][2147483647,1]A.2}{[1][33][]33}}
\newlabel{sol:eje:3-embeddings}{{A.3}{34}{Tema 3}{section*.15}{}}
\newlabel{sol:eje:3-embeddings@cref}{{[subappendix][3][2147483647,1]A.3}{[1][33][]34}}
\newlabel{sol:eje:3-signos}{{A.3}{34}{Tema 3}{section*.16}{}}
\newlabel{sol:eje:3-signos@cref}{{[subappendix][3][2147483647,1]A.3}{[1][34][]34}}
\newlabel{sol:eje:3-capa}{{A.3}{34}{Tema 3}{section*.17}{}}
\newlabel{sol:eje:3-capa@cref}{{[subappendix][3][2147483647,1]A.3}{[1][34][]34}}
\newlabel{sol:eje:3-softmax}{{A.3}{35}{Tema 3}{section*.18}{}}
\newlabel{sol:eje:3-softmax@cref}{{[subappendix][3][2147483647,1]A.3}{[1][35][]35}}
\newlabel{sol:eje:3-nucleo}{{A.3}{35}{Tema 3}{section*.19}{}}
\newlabel{sol:eje:3-nucleo@cref}{{[subappendix][3][2147483647,1]A.3}{[1][35][]35}}
\newlabel{sol:eje:3-word2vec}{{A.3}{35}{Tema 3}{section*.20}{}}
\newlabel{sol:eje:3-word2vec@cref}{{[subappendix][3][2147483647,1]A.3}{[1][35][]35}}
\newlabel{sol:eje:4-vf}{{A.4}{36}{Tema 4}{section*.21}{}}
\newlabel{sol:eje:4-vf@cref}{{[subappendix][4][2147483647,1]A.4}{[1][36][]36}}
\newlabel{sol:eje:4-taylor}{{A.4}{36}{Tema 4}{section*.22}{}}
\newlabel{sol:eje:4-taylor@cref}{{[subappendix][4][2147483647,1]A.4}{[1][36][]36}}
\newlabel{sol:eje:4-mit}{{A.4}{37}{Tema 4}{section*.23}{}}
\newlabel{sol:eje:4-mit@cref}{{[subappendix][4][2147483647,1]A.4}{[1][37][]37}}
\newlabel{sol:eje:4-tasa}{{A.4}{37}{Tema 4}{section*.24}{}}
\newlabel{sol:eje:4-tasa@cref}{{[subappendix][4][2147483647,1]A.4}{[1][37][]37}}
\newlabel{sol:eje:4-necesaria}{{A.4}{38}{Tema 4}{section*.25}{}}
\newlabel{sol:eje:4-necesaria@cref}{{[subappendix][4][2147483647,1]A.4}{[1][38][]38}}
\newlabel{sol:eje:5-vf}{{A.5}{38}{Tema 5}{section*.26}{}}
\newlabel{sol:eje:5-vf@cref}{{[subappendix][5][2147483647,1]A.5}{[1][38][]38}}
\newlabel{sol:eje:5-hessianos}{{A.5}{38}{Tema 5}{section*.27}{}}
\newlabel{sol:eje:5-hessianos@cref}{{[subappendix][5][2147483647,1]A.5}{[1][38][]38}}
\newlabel{sol:eje:5-boyd}{{A.5}{39}{Tema 5}{section*.28}{}}
\newlabel{sol:eje:5-boyd@cref}{{[subappendix][5][2147483647,1]A.5}{[1][39][]39}}
\newlabel{sol:eje:5-red}{{A.5}{39}{Tema 5}{section*.29}{}}
\newlabel{sol:eje:5-red@cref}{{[subappendix][5][2147483647,1]A.5}{[1][39][]39}}
\newlabel{sol:eje:5-softmax-W}{{A.5}{40}{Tema 5}{section*.30}{}}
\newlabel{sol:eje:5-softmax-W@cref}{{[subappendix][5][2147483647,1]A.5}{[1][40][]40}}
\newlabel{sol:eje:5-unico}{{A.5}{40}{Tema 5}{section*.31}{}}
\newlabel{sol:eje:5-unico@cref}{{[subappendix][5][2147483647,1]A.5}{[1][40][]40}}
\newlabel{sol:eje:6-porque}{{A.6}{41}{Tema 6}{section*.32}{}}
\newlabel{sol:eje:6-porque@cref}{{[subappendix][6][2147483647,1]A.6}{[1][41][]41}}
\newlabel{sol:eje:6-cubo}{{A.6}{41}{Tema 6}{section*.33}{}}
\newlabel{sol:eje:6-cubo@cref}{{[subappendix][6][2147483647,1]A.6}{[1][41][]41}}
\newlabel{sol:eje:6-direccional}{{A.6}{42}{Tema 6}{section*.34}{}}
\newlabel{sol:eje:6-direccional@cref}{{[subappendix][6][2147483647,1]A.6}{[1][42][]42}}
\newlabel{sol:eje:6-cs231n}{{A.6}{42}{Tema 6}{section*.35}{}}
\newlabel{sol:eje:6-cs231n@cref}{{[subappendix][6][2147483647,1]A.6}{[1][42][]42}}
\newlabel{sol:eje:6-fermi}{{A.6}{42}{Tema 6}{section*.36}{}}
\newlabel{sol:eje:6-fermi@cref}{{[subappendix][6][2147483647,1]A.6}{[1][42][]42}}
\newlabel{sol:eje:6-hoptimo}{{A.6}{43}{Tema 6}{section*.37}{}}
\newlabel{sol:eje:6-hoptimo@cref}{{[subappendix][6][2147483647,1]A.6}{[1][43][]43}}
\makeatother
